\documentclass[11pt]{article}
\usepackage[margin=1in]{geometry}
\usepackage{amsmath,amssymb}
\usepackage{booktabs}
\usepackage{longtable}
\usepackage{enumitem}
\usepackage{hyperref}
\hypersetup{colorlinks=true, linkcolor=blue, urlcolor=blue, citecolor=blue}
\usepackage{xcolor}
\usepackage{tcolorbox}
\tcbuselibrary{breakable}

\newtcolorbox{warningbox}{colback=red!5!white, colframe=red!60!black, breakable, title=Implementation Warning}
\newtcolorbox{decisionbox}{colback=blue!5!white, colframe=blue!50!black, breakable, title=Design Decision}

\title{A Linearized Multi-Theme CGE Approach to a Real-Time \\
Sectoral--Regional Layoff-Risk Dashboard for Indonesia: \\
Methodology and Data Retrieval Pipeline}
\author{}
\date{\today}

\begin{document}
\maketitle
\tableofcontents
\newpage

\section{Purpose and Overview}

This document specifies the methodology and complete data-retrieval pipeline
for a dashboard that maps current economic conditions into a
\textbf{52-sector $\times$ 34-province} matrix of employment change, used as
an early-warning layoff-risk signal for Indonesia. The dashboard is intended
for use by Indonesia's Ministry of Manpower (Kemnaker) and is built on
IndoTERM, a computable general equilibrium (CGE) model of the Indonesian
economy.

\subsection{The Core Design Problem}

A full CGE solve is computationally expensive and cannot be run live every
time a user presses ``update.'' The design solves this by separating two
steps with very different cost profiles:

\begin{enumerate}
  \item \textbf{Offline calibration} (expensive, infrequent): IndoTERM is
  solved once per shock, at a small (1\%) perturbation, to obtain a
  calibrated elasticity --- the percentage employment response in every
  sector--province cell to a unit shock. This step is run by the user
  outside the dashboard and only needs to be repeated when the model,
  closure rules, or base-year Social Accounting Matrix (SAM) change.
  \item \textbf{Live aggregation} (cheap, every refresh): on each dashboard
  update, the system retrieves the latest values of various real-world
  economic indicators, converts them into shock magnitudes, and multiplies
  them against the pre-calibrated elasticities --- a simple tensor
  contraction, not a CGE solve.
\end{enumerate}

This is a deliberate, documented shortcut, not an oversight. Section~\ref{sec:limitations}
states explicitly which assumptions this rests on, and Section~\ref{sec:future}
describes the planned transition away from this shortcut once server-side,
on-demand CGE solving becomes feasible.

\subsection{Three Shock Themes}

Three conceptually distinct families of shocks feed into the same
sector--province matrix:

\begin{description}
  \item[Theme 1 --- Commodity Price Shocks.] World/domestic price changes in
  eight key Indonesian export and import commodities, directly affecting
  the sectors that produce or use them.
  \item[Theme 2 --- Trading-Partner Growth Shocks.] Revisions to IMF growth
  forecasts for Indonesia's trading partners, translated into
  export-demand shocks for tradable sectors via each sector's actual
  country-by-country export exposure.
  \item[Theme 3 --- Domestic Demand Shocks.] Deviations of domestic retail
  sales volume from trend, by category, capturing slowdowns in
  domestically-oriented sectors that the first two themes, being
  externally focused, do not.
\end{description}

Each theme constructs its own shock magnitudes from different real-world
data sources (detailed in Sections~\ref{sec:theme1}--\ref{sec:theme3}) and is
multiplied against its own, separately-calibrated elasticity tensor (a
deliberate modeling choice explained in Section~\ref{sec:elasticity-scope}).
The three themes' contributions are then summed (Section~\ref{sec:aggregation})
into a single severity-weighted heatmap (Section~\ref{sec:heatmap}).

\section{Notation}

\begin{longtable}{@{}p{3cm}p{11.5cm}@{}}
\toprule
Symbol & Meaning \\
\midrule
\endhead
$i$ & Sector index, $i = 1,\dots,52$ (IndoTERM's \texttt{AGGIND} aggregation) \\
$r$ & Province index, $r = 1,\dots,34$ \\
$j$ & Generic shock index within a theme \\
$\theta$ & Theme index, $\theta \in \{1,2,3\}$ \\
$\eta_{i,r,j}^{(\theta)}$ & Calibrated elasticity: \% change in employment in
  cell $(i,r)$ induced by a 1\% shock to driver $j$ under theme $\theta$,
  holding all other shocks at zero \\
$x_j$ & Live shock magnitude for driver $j$, the dashboard's run-time input \\
$E_{i,r}^{(\theta)}$ & Theme $\theta$'s contribution to total \% employment
  change in cell $(i,r)$ \\
$E_{i,r}$ & Total \% employment change in cell $(i,r)$, summed across all
  three themes \\
$L_{i,r}^{0}$ & Baseline (most recently observed) employment level in cell
  $(i,r)$ \\
$\Delta L_{i,r}$ & Absolute employment change in cell $(i,r)$ --- the
  dashboard's primary heatmap output \\
$\mathcal{K}$ & The subset of the 52 sectors with material export exposure
  (Theme 2's shock-target set) \\
$d^g_c$ & Revision to country $c$'s IMF growth forecast between consecutive
  WEO vintages, for the same target year \\
$w_{k,c}$ & Country $c$'s share of sector $k$'s total Indonesian exports \\
$g_{i,t}^{yoy}$ & Year-on-year growth rate of real retail sales (IPR) for
  category $i$ at month $t$ \\
$\bar g_{i,t}$ & Trailing 36-month moving average of $g_{i,t}^{yoy}$ \\
\bottomrule
\end{longtable}

\section{Sector and Region Classification}
\label{sec:classification}

\subsection{The 52-Sector Aggregation}

Sectors follow IndoTERM's own 52-industry aggregation (\texttt{AGGIND} /
\texttt{AGGCOM}), itself an aggregation of a 185-sector base classification
(\texttt{IND} / \texttt{COM}), with the many-to-one mapping recorded in
\texttt{MIND} / \texttt{MCOM}. The source file (\texttt{sec52.agg}) is a
GEMPACK HAR (Header Array) file --- a binary format requiring GEMPACK or the
\texttt{HARpy} Python library to parse programmatically; it is \emph{not} a
plain-text or CSV file, and generic text/CSV readers will not parse it
correctly.

The 52 sectors, in their native order, are listed in
Table~\ref{tab:sector52}.

\begin{longtable}{@{}p{0.6cm}p{2.8cm}p{0.6cm}p{2.8cm}p{0.6cm}p{2.8cm}p{0.6cm}p{2.8cm}@{}}
\caption{The 52-sector IndoTERM aggregation.}
\label{tab:sector52} \\
\toprule
\# & Sector & \# & Sector & \# & Sector & \# & Sector \\
\midrule
\endhead
1 & FoodCrops & 14 & Textiles & 27 & OtherMan & 40 & Hotels \\
2 & HortiCrops & 15 & Leather & 28 & Electricity & 41 & Restaurant \\
3 & Estates & 16 & WoodProd & 29 & CityGas & 42 & Telecom \\
4 & Livestock & 17 & PaperProd & 30 & WasteMan & 43 & Finance \\
5 & AgricService & 18 & NonMetalProd & 31 & Construction & 44 & Insurance \\
6 & Forestry & 19 & CoalOilMan & 32 & VehicTrade & 45 & OthFinance \\
7 & Fishery & 20 & Chemical & 33 & OthTrade & 46 & FinanSvc \\
8 & Coal & 21 & Rubber & 34 & RailTransp & 47 & RealEstate \\
9 & OilGasGeo & 22 & BasicMetal & 35 & LandTransp & 48 & BusSvc \\
10 & IronOre & 23 & MetalProd & 36 & SeaTransp & 49 & GovAdmin \\
11 & OtherMine & 24 & Machinery & 37 & WaterTransp & 50 & Education \\
12 & FoodMan & 25 & TranspEquip & 38 & AirTransp & 51 & HealthSvc \\
13 & Tobacco & 26 & Furniture & 39 & TranspSvc & 52 & OtherSvc \\
\bottomrule
\end{longtable}

\noindent Codes 1--7 correspond to agriculture/forestry/fishery, 8--11 to
mining, 12--27 to manufacturing (16 industries --- by far the largest single
bucket, and where most of Theme 1's commodities and most of Theme 3's
mapping decisions concentrate), and 28--52 to utilities, construction, and
services (25 industries, almost entirely non-tradable and outside Theme 2's
shock-target set $\mathcal{K}$).

\subsection{Provinces}

Regions correspond to Indonesia's 34 provinces, matching the dimension used
in BPS's Inter-Regional Input-Output Table (52 Industries $\times$ 34
Provinces, 2016), which also supplies the baseline employment/output data
referenced in Section~\ref{sec:heatmap}.

\section{Theme 1: Commodity Price Shocks}
\label{sec:theme1}

\subsection{Mechanism}

Eight key commodities, chosen for their combination of (a) large weight in
Indonesia's export or import structure, and (b) geographic concentration
making them informative at the province level, are each calibrated as an
independent 1\% world-price perturbation in IndoTERM.

\subsection{Commodity List and Data Sources}

\begin{longtable}{@{}p{2.3cm}p{3.6cm}p{4cm}p{3.5cm}@{}}
\caption{Theme 1 commodities, reference series, and data access.}
\label{tab:theme1-sources} \\
\toprule
Commodity & Reference series & Source & Access \\
\midrule
\endhead
CPO / palm oil & Bursa Malaysia FCPO & Commodities-API; or World Bank Pink
  Sheet & Paid REST API (real-time) or free monthly XLSX \\
Coal & Newcastle, 6000 kcal/kg & World Bank Pink Sheet & Free monthly XLSX \\
Nickel & LME cathodes & World Bank Pink Sheet & Free monthly XLSX \\
Copper & LME grade A & World Bank Pink Sheet & Free monthly XLSX \\
TGF (textiles/garments/footwear) & demand-side; no single world price &
  see Theme 2 for the demand-side construction & --- \\
Electronics & demand-side; no single world price & PHLX Semiconductor Index
  (SOX); or BLS export price index, FRED series \texttt{IY3344} & Market
  data API (real-time); or FRED API (free, monthly) \\
Rubber & Asia TSR 20, SICOM nearby & World Bank Pink Sheet & Free monthly
  XLSX \\
Oil/gas & Brent/Dubai/WTI average & World Bank Pink Sheet (monthly); or
  FRED/EIA (daily) & Free API \\
\bottomrule
\end{longtable}

\noindent The World Bank Pink Sheet (\texttt{CMO-Historical-Data-Monthly.xlsx},
fixed URL, updated monthly, no API key) is the single source covering coal,
nickel, copper, rubber, and oil/gas --- five of the eight commodities --- and
is the recommended primary source for those five. Oil/gas additionally has a
genuine daily-frequency free REST API (FRED or EIA) if higher refresh
frequency than monthly is desired for that one commodity specifically.

\subsection{Shock Magnitude Construction}

On each refresh, the latest available value of the raw driver, $X_{j,t}$,
and its value one period prior, $X_{j,t-1}$, are retrieved, and:
\begin{equation}
  x_j = 100 \times \frac{X_{j,t} - X_{j,t-1}}{X_{j,t-1}}.
  \label{eq:theme1-xj}
\end{equation}

\begin{warningbox}
The formula $x_j = 100 \cdot X_{j,t}/X_{j,t-1}$, without subtracting 1
(equivalently, without subtracting $X_{j,t-1}$ in the numerator), does
\emph{not} return zero when $X_{j,t} = X_{j,t-1}$, and will silently
corrupt every downstream calculation. \textbf{Enforce with a pipeline unit
test:} feed a synthetic flat series (no period-over-period change) into the
pipeline and assert that $x_j = 0$ exactly. If $t$ and $t-1$ are not
calendar years, the chosen frequency must match the frequency used in the
elasticity calibration (Section~\ref{sec:theme1-calibration}), and every
series should carry an explicit timestamp/frequency tag so mismatched
vintages across commodities are visible rather than silently blended.
\end{warningbox}

\subsection{Elasticity Calibration}
\label{sec:theme1-calibration}

For each of the 8 commodities, IndoTERM is solved once with a 1\%
perturbation applied to that commodity's price alone, all other shocks held
at zero:
\begin{equation}
  \eta_{i,r,j}^{(1)} = e_{i,r,j}\Big|_{x_j = 1\%}, \qquad j = 1,\dots,8,
  \label{eq:theme1-eta}
\end{equation}
producing \textbf{8 full $52\times34$ elasticity matrices}. This step is run
offline by the user in IndoTERM and is not automated by the data-retrieval
pipeline described in this document.

\section{Theme 2: Trading-Partner Growth Shocks}
\label{sec:theme2}

\subsection{Mechanism}

A trading partner's growth-forecast revision is treated as a uniform,
economy-wide demand shifter facing \emph{all} Indonesian exports to that
country: a 1 percentage-point (pp) revision to country $c$'s growth forecast
is assumed to produce a 1\% change in demand for any Indonesian export to
$c$ (unit pass-through). Product-level differentiation in the resulting
shock comes entirely from \emph{which countries buy how much of each
product} --- i.e., from each sector's own export destination profile --- not
from a separately estimated income elasticity per product or per country.

\begin{decisionbox}
This unit-pass-through assumption is a deliberate simplification, not an
estimated parameter. Real income elasticities of import demand are often
estimated above 1 (sometimes well above, for industrial/commodity inputs
such as nickel or steel). Estimating product/country-specific elasticities
from historical Comtrade export values regressed against historical partner
GDP is the natural upgrade path (Section~\ref{sec:future}).
\end{decisionbox}

\subsection{Step 1: Country Growth Revisions}

For each country $c$ covered by the IMF World Economic Outlook (WEO), the
forecast revision for the same target year $t$, comparing the current
release to the immediately preceding one, is:
\begin{equation}
  d^g_c = g_{c,t}^{\text{(current vintage)}} - g_{c,t}^{\text{(prior vintage)}}.
  \label{eq:theme2-dgc}
\end{equation}

\paragraph{Data source.} IMF World Economic Outlook database, accessed via
the IMF's free SDMX 2.1/3.0 REST API (base URL \texttt{api.imf.org/external/sdmx/3.0},
dataflow \texttt{WEO}, no API key required). Example query construction (Python,
via the \texttt{sdmx1} library or direct REST calls):
\begin{verbatim}
import sdmx
IMF_DATA = sdmx.Client('IMF_DATA')
data_msg = IMF_DATA.data('WEO', key='CHN.NGDP_RPCH', params={'startPeriod': '2024'})
\end{verbatim}
where \texttt{NGDP\_RPCH} is the WEO indicator code for real GDP growth and
\texttt{CHN} is the ISO country code (China).

\begin{warningbox}
WEO is published only \textbf{3 times per year} (April, September/October,
and a January interim Update). The dashboard's ``press update'' refresh
\textbf{will not see a new value for $d^g_c$ between these three release
dates} --- this is a genuine constraint of the underlying data, not a
limitation of the pipeline. The dashboard UI should indicate the WEO vintage
currently in use (e.g., ``as of April 2026 WEO'') rather than imply daily
freshness for this theme.
\end{warningbox}

\subsection{Step 2: Export-Share Weights}

For each tradable sector $k \in \mathcal{K}$ (defined in
Section~\ref{sec:theme2-K}), Indonesia's export value to its destination
countries is retrieved and converted into shares:
\begin{equation}
  w_{k,c} = \frac{X_{k,c}}{\sum_{c'} X_{k,c'}}.
  \label{eq:theme2-weights}
\end{equation}

\paragraph{Data sources (either is viable; both are free and key-based).}
\begin{itemize}
  \item \textbf{UN Comtrade}: free API key registration; up to 100,000
  records per query and 500 API calls per day on the free tier. Official
  Python wrapper \texttt{comtradeapicall} or R wrapper \texttt{comtradr}
  simplify querying by reporter (Indonesia), partner country, HS code, and
  flow direction (exports).
  \item \textbf{BPS WebAPI} (\texttt{webapi.bps.go.id}): free API key; the
  \texttt{dataexim} resource exposes Indonesia's own export/import data
  directly by HS code, in JSON, e.g.
  \begin{verbatim}
https://webapi.bps.go.id/v1/api/dataexim/sumber/{sumber}/kodehs/{kodehs}
/jenishs/{jenishs}/tahun/{tahun}/periode/{periode}/key/{key}
  \end{verbatim}
  with an official Python wrapper (\texttt{stadata}) available.
\end{itemize}

\begin{decisionbox}
Export-share weights $w_{k,c}$ change slowly (annual trade patterns, not
month-to-month), so this step does \textbf{not} need to be re-run on every
dashboard refresh. Recompute annually (or whenever a new year of Comtrade/BPS
trade data becomes available) and cache; only $d^g_c$ (Step 1) needs to be
checked against the live WEO release schedule on every refresh.
\end{decisionbox}

\subsection{Step 3: Aggregation into a Sector-Specific Shock}

The shock magnitude for sector $k$ is the export-share-weighted average of
the growth-forecast revisions facing its actual buyers:
\begin{equation}
  x_k = q_k^T = \sum_{c} w_{k,c} \cdot d^g_c.
  \label{eq:theme2-agg}
\end{equation}

\paragraph{Derivation note.} An earlier, more general formulation considered
a separate income elasticity $\epsilon_{i,c}$ per (product, country) pair,
giving $q_{i,c} = \epsilon_i \cdot d^g_c$ aggregated as
$q_i^T = \sum_c w_{i,c}\epsilon_i d^g_c = \epsilon_i \sum_c w_{i,c} d^g_c$.
Under the unit-pass-through assumption ($\epsilon_i \equiv 1$ for all
products, applied uniformly across countries rather than varying by
country), this collapses exactly to Equation~\eqref{eq:theme2-agg}: there is
no need to ever materialize a per-country shock $q_{k,c}$ individually, only
the weighted average.

\subsection{The Tradable Sector Set $\mathcal{K}$}
\label{sec:theme2-K}

Not all 52 sectors have meaningful export exposure; shocking sectors with
near-zero export shares would add CGE-calibration cost for negligible
benefit. $\mathcal{K}$ is determined empirically from the export column of
the 52-sector Input-Output table (export value $\div$ total output per
sector, against a chosen threshold --- 1\% or 5\% are reasonable candidates;
the exact cutoff is a user decision not yet finalized).

\begin{warningbox}
The sector list below is a \textbf{category-level estimate} based on which
of the 52 sectors are agriculture/mining/manufacturing (vs.\ services,
utilities, construction, which are excluded with high confidence). It has
not yet been verified row-by-row against actual export/output ratios in the
IO table. Treat as a starting point for verification, not a final list.
\end{warningbox}

Candidate $\mathcal{K}$ (by sector number and name, from
Table~\ref{tab:sector52}): 3 (Estates), 7 (Fishery), 8 (Coal), 9 (OilGasGeo),
10 (IronOre), 11 (OtherMine), 12 (FoodMan), 14 (Textiles), 15 (Leather), 16
(WoodProd), 19 (CoalOilMan), 20 (Chemical), 21 (Rubber), 22 (BasicMetal), 23
(MetalProd), 24 (Machinery), 25 (TranspEquip), 26 (Furniture) --- approximately
\textbf{17--18 sectors}.

\subsection{Elasticity Calibration}

For each sector $k \in \mathcal{K}$, IndoTERM is solved once with a 1\%
export-demand perturbation to sector $k$ alone:
\begin{equation}
  \eta_{i,r,k}^{(2)} = e_{i,r,k}\Big|_{x_k = 1\%}, \qquad k \in \mathcal{K},
\end{equation}
producing \textbf{approximately 17 full $52\times34$ elasticity matrices}.

\section{Theme 3: Domestic Demand Shocks}
\label{sec:theme3}

\subsection{Mechanism and Motivation}

Themes 1 and 2 are both externally focused (world prices, foreign growth)
and so are silent on purely domestic demand conditions --- relevant for
sectors selling mainly to Indonesian consumers (retail, domestically-oriented
manufacturing, food processing for the domestic market). Theme 3 fills this
gap using Bank Indonesia's real (volume) retail sales index.

\subsection{Why a Real (Volume) Index, Not a Price Index}

An earlier design considered Indonesia's Consumer Price Index (CPI) as the
domestic-demand signal. This was rejected: a price index alone cannot
distinguish a demand-driven price change from a supply/cost-push price
change --- a price rise is consistent with either rising demand (price up,
quantity up) or a negative supply shock (price up, quantity down), and these
have opposite implications for layoff risk. The originally proposed
resolution was a price--quantity co-movement check (pairing CPI with a
volume series and only treating a price move as demand-driven when price and
quantity move in the same direction). This check was \textbf{dropped} once
the design moved to Bank Indonesia's \emph{Indeks Penjualan Riil} (IPR,
``Real Sales Index''): IPR is already a deflated, real/volume measure, so
the price/demand identification problem that motivated the co-movement check
does not arise from this series in the first place, and the extra CPI-fetch
step would be solving an already-solved problem.

\begin{warningbox}
\textbf{Accepted, transparent simplification:} this reasoning assumes Bank
Indonesia's internal deflation of IPR is adequate. If BI's deflator
imperfectly strips out category-specific price effects (e.g.\ uses a
general rather than category-specific price index internally), some
residual price-driven variation could still leak into the yoy growth figure
used as Theme 3's shock. This is not independently re-verified against CPI
in this version of the design --- a contestable but openly stated choice,
not an oversight.
\end{warningbox}

\subsection{Data Source}

Bank Indonesia's Survei Penjualan Eceran (SPE, Retail Sales Survey),
specifically its headline output, the Indeks Penjualan Riil (IPR). Published
monthly. There is no key-based REST API; data is distributed as a
downloadable zip file attached to each month's release page, at a
predictable URL pattern:
\begin{verbatim}
https://www.bi.go.id/id/publikasi/laporan/Documents/
Data-Series-SPE-{Month}-{Year}.zip
\end{verbatim}
where \texttt{\{Month\}} is the Indonesian month name (e.g.\ \texttt{Februari})
and \texttt{\{Year\}} the four-digit year. The release page itself follows the
analogous pattern \texttt{.../Pages/SPE-\{Month\}-\{Year\}.aspx}.

\subsection{``Press Update'' Logic for Theme 3}

\begin{enumerate}
  \item Construct the expected URL for the current month/year using the
  pattern above.
  \item Issue a HEAD request (or attempt a fetch with graceful fallback) to
  check whether this URL exists and differs from the last cached vintage.
  \item If a newer month's zip is available: download, unzip, extract the
  category-level yoy growth series, append to the cached historical panel.
  \item If not (i.e.\ the current month's release has not yet been
  published): proceed using the most recently cached vintage; do not error.
\end{enumerate}

\subsection{IPR Categories}

IPR publishes seven categories, retained in full in the data pipeline
regardless of whether each one is ultimately wired into an elasticity (see
Section~\ref{sec:theme3-mapping}):

\begin{enumerate}
  \item Makanan, Minuman, dan Tembakau (Food, Beverages, Tobacco)
  \item Sandang (Clothing/Textiles)
  \item Suku Cadang dan Aksesori (Automotive Spare Parts and Accessories)
  \item Bahan Bakar Kendaraan Bermotor (Motor Vehicle Fuel)
  \item Barang Budaya dan Rekreasi (Cultural and Recreational Goods)
  \item Peralatan Informasi dan Komunikasi (Information and Communication
  Equipment)
  \item Perlengkapan Rumah Tangga Lainnya (Other Household Equipment)
\end{enumerate}

\subsection{Why Year-on-Year, Not Month-to-Month}

IPR publishes growth both year-on-year (yoy) and month-to-month (mtm).
\textbf{Only the yoy figure is used.} Indonesian retail sales carry strong,
recurring seasonality tied to the national religious/holiday calendar
(Ramadan, Idulfitri, Christmas/New Year), and these dates shift by
approximately 11 days earlier each Gregorian year for the lunar-calendar
holidays. Comparing month $t$ to month $t-12$ (yoy) automatically cancels
this seasonality; compounding mtm growth into an annualized rate does not,
and will badly overstate or understate the true signal whenever a holiday
period falls in the observation window. (BPS/BI press releases routinely and
explicitly attribute large mtm swings to specific holiday periods,
confirming the size of this seasonal effect.)

\subsection{Benchmark/Trend and Shock Construction}

Raw yoy growth is always technically signed correctly (it can be negative),
but a positive-yet-low growth rate still represents a meaningful
\emph{slowdown} relative to a sector's normal trend, and a layoff-risk
signal needs to register that as a negative shock --- not as ``still
growing, so still fine.'' The shock is therefore constructed as a
\textbf{deviation from a trailing benchmark}, not the raw yoy growth rate
itself:
\begin{equation}
  \bar g_{i,t} = \frac{1}{36}\sum_{k=1}^{36} g_{i,t-k}^{yoy},
  \qquad
  x_i = g_{i,t}^{yoy} - \bar g_{i,t}.
  \label{eq:theme3-shock}
\end{equation}

\begin{decisionbox}
A trailing 36-month moving average was chosen over (a) a fixed historical
base-period mean, which requires an arbitrary choice of reference window and
does not adapt to genuine structural shifts (e.g.\ permanent changes in
e-commerce penetration affecting the ICT-equipment category), and (b) formal
trend extraction (HP filter or similar), judged disproportionate for a
first-version dashboard and subject to end-point bias exactly where the
signal is most needed. The 36-month window is long enough to average out
holiday-timing noise across different years while remaining short enough to
track genuine multi-year drift.
\end{decisionbox}

\begin{warningbox}
$\bar g_{i,t}$ requires \textbf{at least 36 months of historical data before
it can be computed for the first time}. The initial pipeline setup must
backfill several years of past SPE zip files (not just the current month)
before Theme 3 can produce any output. Each subsequent monthly refresh only
needs to add one new observation and slide the window forward.
\end{warningbox}

\subsection{Regional Dimension: National Shock, Regional Elasticity}
\label{sec:theme3-regional}

IPR's underlying survey covers a limited set of cities, not all 34
provinces, so $x_i$ is constructed and used as a \textbf{single national
value per category} --- there is no $x_{i,r}$ varying by province. All
regional differentiation in Theme 3's contribution to the heatmap comes
entirely from the elasticity tensor $\eta^{(3)}_{i,r,m}$, which \emph{is}
fully calibrated at $52\times34$ resolution. This means a national
slowdown in, say, domestic clothing sales is assumed to affect every
province's clothing-sector employment in the same relative direction, with
only the magnitude of that effect varying by province (through $\eta$) ---
genuine sub-national divergence in domestic demand itself (one province
softening while another strengthens) is not captured by this version.

\begin{decisionbox}
This national-shock/regional-elasticity structure is not unique to Theme 3
--- Theme 1 (a world commodity price) and Theme 2 (a country's growth
forecast) are also necessarily national-level shocks, since neither has an
Indonesian regional dimension to begin with. All three themes therefore
share the same underlying structural pattern: a single national shock
magnitude, multiplied against a fully region-differentiated elasticity
tensor.
\end{decisionbox}

\subsection{Mapping IPR Categories to the 52-Sector Classification}
\label{sec:theme3-mapping}

This mapping is performed manually by the model's developer, not
automated by the data pipeline, since several IPR categories are diffuse
enough to require judgment (and some may legitimately split across more
than one of the 52 sectors). Indicative mappings, based on category content
alone:

\begin{itemize}
  \item Makanan, Minuman, dan Tembakau $\to$ FoodMan (12), possibly
  Tobacco (13)
  \item Sandang $\to$ Textiles (14), possibly Leather (15)
  \item Suku Cadang dan Aksesori $\to$ TranspEquip (25)
  \item Bahan Bakar Kendaraan Bermotor $\to$ OilGasGeo (9) and/or
  CoalOilMan (19)
  \item Barang Budaya dan Rekreasi $\to$ diffuse; candidate sectors include
  OtherMan (27), OthTrade (33), or a services sector, pending the
  developer's judgment
  \item Peralatan Informasi dan Komunikasi $\to$ Machinery (24) (where
  electronics is aggregated in this classification)
  \item Perlengkapan Rumah Tangga Lainnya $\to$ Furniture (26), possibly
  splitting with other manufacturing sectors
\end{itemize}

\subsection{Elasticity Calibration: All 52 Sectors}
\label{sec:theme3-calibration}

\begin{decisionbox}
Unlike Theme 2, where the elasticity-target set was restricted to
$\mathcal{K}$ (tradable sectors only) because non-tradable sectors have
negligible export-demand exposure by construction, Theme 3's elasticity
target set is \textbf{all 52 sectors, with no pre-filtering by IPR's current
category coverage}. The reasoning: every sector sells \emph{something}
domestically, even sectors IPR's seven categories do not currently observe;
whether a sector's domestic-demand elasticity turns out to be economically
negligible is itself a CGE result to be discovered, not an assumption to be
imposed in advance based on today's data availability. Calibrating all 52
up front means that if more granular domestic-demand data becomes available
later (e.g.\ from Susenas or an expanded retail survey), the elasticity side
of the pipeline is already complete and does not block using it.
\end{decisionbox}

For each of the 52 sectors, IndoTERM is solved once with a 1\%
domestic-demand perturbation to that sector alone:
\begin{equation}
  \eta_{i,r,m}^{(3)} = e_{i,r,m}\Big|_{x_m = 1\%}, \qquad m = 1,\dots,52,
\end{equation}
producing \textbf{52 full $52\times34$ elasticity matrices} --- the largest
single calibration task among the three themes.

\section{Why Elasticities Are Theme-Specific}
\label{sec:elasticity-scope}

A central design decision, worth stating explicitly because it substantially
affects the total CGE-calibration workload: \textbf{a sector's elasticity
tensor is not shared across themes}, even when the same sector is the
directly-shocked target under more than one theme (e.g.\ Textiles appears in
Theme 1 as an export-price target and, via the Sandang mapping, in Theme 3
as a domestic-demand target).

The justification is economic, not merely a bookkeeping convention: a
domestic-demand shock and an export-price shock to the same sector are
mechanically different perturbations to IndoTERM. They enter the model
through different channels (domestic final consumption vs.\ export
price/demand), and they can plausibly activate different substitution
elasticities and different general-equilibrium feedback paths as they
propagate through the input-output structure --- domestic consumers and
export markets need not substitute toward or away from a sector's output in
the same way. Treating $\eta^{(1)}_{i,r,\text{Textiles}}$ and
$\eta^{(3)}_{i,r,\text{Textiles}}$ as identical would be assuming away a real
distinction rather than testing or justifying it.

The consequence is a larger total calibration workload than a
shared-elasticity design would require, summarized in
Table~\ref{tab:elasticity-scope}.

\begin{table}[h]
\centering
\begin{tabular}{@{}lccl@{}}
\toprule
Theme & \# of shocks & \# of $52\times34$ matrices & Calibration basis \\
\midrule
1 --- Commodity price & 8 & 8 & 1\% price perturbation, 8 commodities \\
2 --- Partner growth & $\approx 17$ ($\mathcal{K}$) & $\approx 17$ & 1\%
  export-demand perturbation, tradable sectors only \\
3 --- Domestic demand & 52 & 52 & 1\% domestic-demand perturbation, all
  sectors \\
\midrule
\textbf{Total} & & $\approx 77$ & \\
\bottomrule
\end{tabular}
\caption{Elasticity-calibration scope across the three themes. This is the
CGE-modeling workload, supplied externally by the model developer; it is
not automated by the data-retrieval pipeline described in this document.}
\label{tab:elasticity-scope}
\end{table}

\section{Aggregation Across Themes}
\label{sec:aggregation}

Each theme's contribution to the percentage employment change in cell
$(i,r)$ is:
\begin{align}
  E_{i,r}^{(1)} &= \sum_{j=1}^{8} \eta_{i,r,j}^{(1)} \cdot x_j,
  \label{eq:E1} \\
  E_{i,r}^{(2)} &= \sum_{k \in \mathcal{K}} \eta_{i,r,k}^{(2)} \cdot q_k^T,
  \qquad q_k^T = \sum_{c} w_{k,c} \cdot d^g_c,
  \label{eq:E2} \\
  E_{i,r}^{(3)} &= \sum_{m=1}^{52} \eta_{i,r,m}^{(3)} \cdot
  \left(g_{m,t}^{yoy} - \bar g_{m,t}\right).
  \label{eq:E3}
\end{align}
The total percentage employment effect is the simple sum across themes:
\begin{equation}
  E_{i,r} = E_{i,r}^{(1)} + E_{i,r}^{(2)} + E_{i,r}^{(3)}.
  \label{eq:Etotal}
\end{equation}

This inherits the additive-separability assumption at two nested levels:
across shocks \emph{within} a theme (Equations~\eqref{eq:E1}--\eqref{eq:E3}),
and across the three \emph{themes} themselves (Equation~\eqref{eq:Etotal}).
Both levels assume no interaction effects between simultaneously active
shocks --- discussed further in Section~\ref{sec:limitations}.

\section{Severity-Weighted Heatmap}
\label{sec:heatmap}

$E_{i,r}$ alone is a \emph{percentage} change, and so weights a large
percentage shock to a small sector--province cell identically to the same
percentage shock hitting a large cell --- the wrong lens for a \emph{risk}
dashboard, which should foreground absolute headcount exposure. The primary
heatmap output is therefore the absolute employment change:
\begin{equation}
  \Delta L_{i,r} = E_{i,r} \cdot L_{i,r}^{0},
  \label{eq:DeltaL}
\end{equation}
where $L_{i,r}^0$ is the baseline (most recently observed) employment level
in cell $(i,r)$, sourced from the same Input-Output table that supplies the
sector and province dimensions (Section~\ref{sec:classification}), or from a
more frequently updated labor-force survey if available --- in which case
the vintage difference between $L_{i,r}^0$'s source and the SAM's base year
should be tagged explicitly in metadata, since using two different-vintage
sources silently is a hidden inconsistency.

$\Delta L_{i,r}$ is rendered as the dashboard's primary heatmap (color
intensity proportional to magnitude, with sign indicating job gains vs.\
losses); $E_{i,r}$ (the pure percentage figure) is retained as a
secondary/tooltip layer, since ``where is the biggest percentage shock'' and
``where is the biggest headcount exposure'' are different and both useful
questions for a policy user to be able to ask.

\section{Computational Architecture: ``Press Update'' Logic}
\label{sec:architecture}

\begin{enumerate}
  \item \textbf{Theme 1.} For each of the 8 commodities, check the latest
  available value from its source (World Bank Pink Sheet monthly file, or
  FRED/EIA daily API for oil/gas) against the last cached value. If newer,
  fetch and recompute $x_j$ via Equation~\eqref{eq:theme1-xj}.
  \item \textbf{Theme 2.} Check whether a new IMF WEO vintage has been
  published since the last cached release (April / September-October /
  January). If so, re-pull $d^g_c$ for all relevant countries via the IMF
  SDMX API and recompute $q_k^T$ for each $k \in \mathcal{K}$ using
  Equation~\eqref{eq:theme2-agg}, combined with the cached export-share
  weights $w_{k,c}$ (these are recomputed on a much slower, e.g.\ annual,
  cadence from Comtrade/BPS trade data --- not on every refresh).
  \item \textbf{Theme 3.} Construct the expected SPE zip URL for the current
  month; check existence/freshness against the cached vintage; if newer,
  download, unzip, append the new month's category-level yoy figures to the
  trailing panel, and recompute $\bar g_{i,t}$ and $x_i$ for all seven IPR
  categories via Equation~\eqref{eq:theme3-shock}.
  \item \textbf{Multiply and aggregate.} Each theme's shock vector is
  multiplied against its respective pre-stored elasticity tensor (a tensor
  contraction, not a CGE solve), summed across themes
  (Equation~\eqref{eq:Etotal}), and scaled by baseline employment
  (Equation~\eqref{eq:DeltaL}).
  \item \textbf{Render.} The resulting $\Delta L_{i,r}$ matrix is rendered as
  the heatmap; $E_{i,r}$ is exposed as supplementary detail.
\end{enumerate}

\begin{decisionbox}
The elasticity tensors themselves (Section~\ref{sec:elasticity-scope}, the
slow and expensive part of the whole pipeline) are \textbf{never recomputed
on refresh}. They are supplied externally by the model developer from
IndoTERM runs and only need to be re-supplied when the CGE model, its
closures, or the base-year SAM change --- not on any dashboard-refresh
cadence.
\end{decisionbox}

\section{Assumptions and Known Limitations}
\label{sec:limitations}

\begin{itemize}
  \item \textbf{Local linearity.} Each $\eta^{(\theta)}_{i,r,j}$ is the slope
  measured at a 1\% perturbation, applied as though constant for any
  observed magnitude of $x_j$. CGE responses are generally nonlinear once
  closures bind differently at different shock sizes --- for instance, the
  dual labor-market closure (rigid-wage/quantity-adjusting formal
  manufacturing vs.\ underemployment-absorbing informal sectors) may not
  respond proportionally as shock size grows. Historical shocks of genuine
  interest (large coal, CPO, or nickel price swings; sizable WEO forecast
  revisions during global downturns) are frequently far outside the small
  neighborhood in which a 1\%-calibrated slope is guaranteed to extrapolate
  cleanly.
  \item \textbf{Additive separability, within and across themes.} Summing
  across shocks within a theme, and summing across the three themes
  (Section~\ref{sec:aggregation}), both assume the absence of interaction
  effects. This holds only approximately under genuine joint
  general-equilibrium market clearing, and the approximation degrades when
  multiple shocks load on the same factor or sector simultaneously --- e.g.\
  Theme 1's TGF export-price shock and Theme 3's Sandang domestic-demand
  shock both affecting Textiles employment in the same period, with no
  cross-term capturing their joint effect.
  \item \textbf{Theme-specific elasticities, no cross-theme reuse} (see
  Section~\ref{sec:elasticity-scope} for full justification). This roughly
  triples the CGE-calibration workload relative to a design that shares one
  elasticity tensor per sector across all themes, but avoids silently
  assuming that domestic-demand and export shocks propagate identically.
  \item \textbf{Theme 3's shock is national, not regional}
  (Section~\ref{sec:theme3-regional}). Genuine sub-national divergence in
  domestic demand --- one province's retail sector softening while
  another's strengthens --- is not captured; only the elasticity side of
  Theme 3 varies by province.
  \item \textbf{IPR's internal deflation is taken at face value.} If Bank
  Indonesia's real-sales deflator imperfectly isolates volume from price
  effects, some residual price-driven noise could leak into Theme 3's shock
  construction. This is an openly stated simplification, not independently
  re-verified against CPI in this version.
  \item \textbf{Unit pass-through in Theme 2.} A 1pp growth-forecast revision
  is assumed to produce exactly a 1\% export-demand change, uniformly across
  products and countries, with no separately estimated income elasticity.
  True elasticities frequently differ from 1, especially for
  industrial/commodity exports.
  \item \textbf{Time-invariant elasticities and baseline employment between
  refreshes.} Both $\eta^{(\theta)}_{i,r,j}$ and $L_{i,r}^0$ are held fixed
  until the user manually recalibrates IndoTERM or re-sources employment
  data; an explicit recalibration cadence for each, and explicit vintage
  tagging where sources differ, should be defined rather than left implicit.
  \item \textbf{Heterogeneous epistemic status of inputs across themes.}
  Theme 1's inputs are CGE-simulated structural elasticities applied to
  directly observed world prices. Theme 2 combines a CGE-simulated
  elasticity with an unestimated unit-pass-through assumption. Theme 3
  combines a CGE-simulated elasticity with an empirical trend-deviation
  measure subject to a deflation-quality caveat. A dashboard user should not
  be given the impression that all three carry the same epistemic weight.
  \item \textbf{Discrete, low-frequency refresh for Theme 2 specifically.}
  Unlike Themes 1 and 3, which can in principle update monthly or even
  daily, Theme 2's $d^g_c$ only changes 3 times per year (WEO release
  schedule) --- the dashboard should communicate this vintage explicitly
  rather than implying continuous freshness.
  \item \textbf{No joint re-solve benchmarking yet.} There is currently no
  periodic check of the full linear-additive approximation
  (Equation~\eqref{eq:Etotal}) against a true joint CGE re-solve under
  realistic, simultaneously active multi-theme shock combinations. Before
  any high-stakes policy communication, this benchmark should be run to
  establish an empirical error bound on the approximation.
\end{itemize}

\section{Future Work}
\label{sec:future}

\begin{itemize}
  \item \textbf{Retire the linear shortcut entirely, rather than patch it.}
  The calibrate-once-and-multiply design (Sections~\ref{sec:theme1}--\ref{sec:architecture})
  is a deliberate stopgap, not the intended end state. Once IndoTERM can be
  run live on a server with acceptable latency, the dashboard should move
  from \emph{look up a pre-calibrated elasticity and multiply} to
  \emph{solve the full CGE on the actual observed combination of shocks on
  every refresh}. This removes the linearity and additive-separability
  assumptions in Section~\ref{sec:limitations} outright, rather than
  approximating around them with interaction terms or quadratic
  corrections. Server sizing, solve time, and a queuing/caching strategy for
  concurrent users are the open questions for that transition, to be scoped
  as a separate piece of work.
  \item \textbf{Two-point calibration check for Theme 1.} For commodities
  with historically large price swings (coal, CPO, nickel), additionally
  calibrate at a 10\% perturbation and compare the implied slope to the
  1\%-calibrated $\eta$. If stable, the linear extrapolation is empirically
  supported for that commodity; if not, fit a quadratic correction term
  using the two calibration points.
  \item \textbf{Estimate product/country-specific elasticities for Theme 2}
  in place of the unit-pass-through assumption, via a log-log regression of
  historical Comtrade export values against historical partner GDP growth,
  by product.
  \item \textbf{Sub-national domestic-demand data for Theme 3,} if/when a
  province-level consumption proxy becomes available (e.g.\ BPS's Susenas
  household expenditure survey) --- though Susenas's much lower publication
  frequency relative to IPR's monthly cadence would require its own
  treatment, not a simple substitution.
  \item \textbf{Interaction terms across themes,} for sectors confirmed to
  be jointly and materially hit by more than one theme simultaneously, once
  the re-solve benchmarking in Section~\ref{sec:limitations} quantifies how
  large the currently-omitted interaction error actually is in practice.
\end{itemize}

\section{Open Items Requiring Further User Input}
\label{sec:open-items}

\begin{itemize}
  \item Verify $\mathcal{K}$ (Section~\ref{sec:theme2-K}) row-by-row against
  actual export/output ratios in the 52-sector Input-Output table; the list
  given is a category-level estimate, not yet confirmed.
  \item Choose the export-share threshold defining $\mathcal{K}$ (1\%? 5\%?
  another value?) --- not yet decided.
  \item Complete the manual mapping of IPR's seven categories onto the
  52-sector classification (Section~\ref{sec:theme3-mapping}); some
  categories may legitimately split across more than one sector.
  \item Specify the source and vintage for $L_{i,r}^0$ (baseline
  employment); not yet finalized, and should be tagged with its own update
  cadence, potentially distinct from the IndoTERM SAM's base year.
\end{itemize}

\end{document}
